\section{Тест производительности}
Тест производительности состоит из трёх частей. В первой результаты жадного алгоритма
сравниваются с наивным перебором на случайных малых тестах, чтобы проверить корректность.
Во второй время жадного алгоритма сравнивается со временем перебора всех подмножеств
добавок при растущем $M$. В третьей время жадного алгоритма измеряется при росте $N$,
чтобы проверить оценку $O(MN^2)$.

Наивный перебор проходит по всем подмножествам добавок, пропускает те, в которых меньше
$N$ добавок или стоимость не меньше лучшей найденной, а для остальных независимо вычисляет ранг
матрицы методом Гаусса с выбором опорного элемента. Среди подмножеств ранга $N$ выбирается
самое дешёвое. Количество подмножеств равно $2^M$, поэтому время перебора растёт экспоненциально.

Время измерялось с помощью \texttt{std::chrono::steady\_clock}. Для каждого алгоритма замер
проводился 3 раза, в таблицах указано среднее
значение. Измерялось только решение (сортировка и метод Гаусса); чтение ввода и вывод не
учитывались, а копирование матрицы входит в замер. Для замеров использовались случайные матрицы
с коэффициентами и ценами от 0 до 50.

\textbf{Корректность.} Результаты жадного алгоритма сравнивались с перебором на
3000 случайных тестах ($N$ от 1 до 5, $M$ от $N$ до $N + 5$, коэффициенты от 0 до 2, цены от 0 до 10).
Малые коэффициенты и цены выбраны специально, чтобы получалось много линейно зависимых строк и
одинаковых цен. В каждом тесте проверялось, что ответ $-1$ получается одновременно у обоих алгоритмов,
что стоимость найденного набора равна оптимальной и что ранг выбранных добавок равен $N$.
Все 3000 тестов пройдены, при этом в 222 из них ответом было $-1$.

\textbf{Сравнение с перебором.} Число веществ $N = 5$, число добавок $M$ растёт.

\begin{table}[h]
\centering
\begin{tabular}{|c|c|c|c|c|}
\hline
$M$ & $N$ & Жадный алгоритм, с & Перебор, с & Отношение \\
\hline
8  & 5 & 0{,}000006880 & 0{,}000028133 & 4{,}1 \\
12 & 5 & 0{,}000007280 & 0{,}000342533 & 47{,}1 \\
16 & 5 & 0{,}000007060 & 0{,}004736333 & 670{,}9 \\
20 & 5 & 0{,}000006400 & 0{,}089583433 & 13\,997{,}4 \\
\hline
\end{tabular}
\end{table}

\textbf{Масштабирование по $N$.} Число добавок $M = 2N$. Измерялись два сценария. В
первом матрица имеет полный ранг, и алгоритм останавливается, как только набран базис.
Во втором последний столбец матрицы нулевой, ранг равен $N - 1$, базис полным не становится,
и алгоритм обрабатывает все $M$ строк, то есть работает в худшем случае.

\begin{table}[h]
\centering
\begin{tabular}{|c|c|c|c|c|c|}
\hline
$N$ & $M$ & Полный ранг, с & Ранг $N-1$, с & Рост & нс/($MN^2$) \\
\hline
50  & 100  & 0{,}000136 & 0{,}000236 & ---  & 0{,}943 \\
100 & 200  & 0{,}000845 & 0{,}001858 & 7{,}88 & 0{,}929 \\
200 & 400  & 0{,}008233 & 0{,}014951 & 8{,}05 & 0{,}934 \\
400 & 800  & 0{,}051325 & 0{,}118204 & 7{,}91 & 0{,}923 \\
800 & 1600 & 0{,}429014 & 1{,}040611 & 8{,}80 & 1{,}016 \\
\hline
\end{tabular}
\end{table}

В столбце <<Рост>> указано отношение времени худшего случая ко времени для предыдущего $N$.
В последнем столбце время худшего случая разделено на $MN^2$ и выражено в наносекундах.

\begin{verbatim}
PS laba-8> .\build\Release\bench.exe
Correctness vs brute force: OK (3000 random tests, 222 of them with answer -1)

M    N    greedy, s        brute force, s   ratio
8    5    0.000006880      0.000028133      4.1
12   5    0.000007280      0.000342533      47.1
16   5    0.000007060      0.004736333      670.9
20   5    0.000006400      0.089583433      13997.4

N      M      full rank, s     rank N-1, s      growth     ns/(M*N^2)
50     100    0.000136         0.000236         -          0.943
100    200    0.000845         0.001858         x7.88      0.929
200    400    0.008233         0.014951         x8.05      0.934
400    800    0.051325         0.118204         x7.91      0.923
800    1600   0.429014         1.040611         x8.80      1.016
\end{verbatim}

Из первой таблицы видно, что время перебора растёт экспоненциально: каждое увеличение
$M$ на 4 умножает его в 12--19 раз, что соответствует росту $2^M$ (при увеличении $M$ на
4 число подмножеств вырастает в 16 раз). При $M = 20$ перебор уже работает около 0{,}09 с и
в 14 тысяч раз медленнее жадного алгоритма, а при больших $M$ он станет неприменим. Время жадного
алгоритма на таких малых данных остаётся на уровне нескольких микросекунд и от $M$ почти не
зависит: оно определяется накладными расходами (выделение памяти, копирование матрицы), поэтому
небольшие колебания не отражают реальной асимптотики.

Вторая таблица подтверждает оценку $O(MN^2)$. Так как $M = 2N$, время должно расти
пропорционально $N^3$, то есть при удвоении $N$ примерно в 8 раз, что и наблюдается
(отношения 7{,}9--8{,}8). Время, приходящееся на единицу $MN^2$, остаётся в пределах
0{,}92--1{,}02~нс, то есть практически постоянно. Небольшое увеличение при $N = 800$
(рост в 8{,}8 раза, 1{,}016~нс) вероятно связано с кэшем процессора: матрица $1600 \times 801$ чисел
типа \texttt{double} занимает около 10~МБ и не помещается в кэш, поэтому обращения к памяти
становятся медленнее. На асимптотику это не влияет. В худшем случае ($N = 800$, $M = 1600$)
решение работает около одной секунды.

Сценарий с полным рангом быстрее худшего случая в 1{,}7--2{,}4 раза, потому что алгоритм
прекращает работу после набора базиса и не обрабатывает оставшиеся строки. Он также
растёт примерно в 8 раз при удвоении $N$.

\pagebreak